\subsubsection{12.01.15}
\begin{enumerate}
	
	\item Время начала и окончания собрания: 17:00 - 21:30.
	
	\item Цели собрания: 
	\begin{enumerate}
		
	    \item Изготовить из листа стали желоб и установить его на робота.
			
		\item Закрепить новый сверхпрочный провод на подъемнике.
		
		\item Проверить насколько точно падают мячи благодаря желобу.
	    
	\end{enumerate}

	\item Проделанная работа:
	\begin{enumerate}
		
		\item Провод был проведен и сервопривод подключен.
		\begin{figure}[H]
			\begin{minipage}[h]{0.2\linewidth}
				\center  
			\end{minipage}
			\begin{minipage}[h]{0.6\linewidth}
				\center{\includegraphics[scale=0.3]{days/12.01.15/images/01}}
				\caption{Провод на подъемнике}
			\end{minipage}
		\end{figure}
		
		\item Сегодня было решено заменить откосы на МЗК с пластиковых на жесткие металлические, которые не будут прогибаться и будут лучше центрировать подвижную корзину.
		\begin{figure}[H]
			\begin{minipage}[h]{0.2\linewidth}
				\center  
			\end{minipage}
			\begin{minipage}[h]{0.6\linewidth}
				\center{\includegraphics[scale=0.19]{days/12.01.15/images/02}}
				\caption{Новые откосы}
			\end{minipage}
		\end{figure}
		
		\item Желоб был изготовлен и установлен на робота.
		\begin{figure}[H]
			\begin{minipage}[h]{0.2\linewidth}
				\center  
			\end{minipage}
			\begin{minipage}[h]{0.6\linewidth}
				\center{\includegraphics[scale=0.2]{days/12.01.15/images/03}}
				\caption{Желоб на роботе}
			\end{minipage}
		\end{figure}
		
		\item К концу желоба была прикреплена пластиковая откидная часть с отверстием, останавливающая мячи и заставляющая их падать отвесно вниз.
		\begin{figure}[H]
			\begin{minipage}[h]{0.47\linewidth}
				\center{\includegraphics[scale=0.2]{days/12.01.15/images/04}}
			\end{minipage}
			\hfill
			\begin{minipage}[h]{0.47\linewidth}
				\center{\includegraphics[scale=0.21]{days/12.01.15/images/05}}
			\end{minipage}
			\caption{Желоб с откидывающейся частью}
		\end{figure}
		
        \item Испытания ковша и желоба с откидной направляющей показали, что мячи выпадают строго вертикально, что обеспечивает потрясающую точность. 
		
		\item Для того, чтобы помещаться в габариты, изначально откидная направляющая находится в сложенном состоянии. Для того, чтобы открыть ее, было решено соединить ее с ковшом леской таким образом, чтобы она открывалась синхронно с опрокидыванием ковша. Было решено, что мы займемся этим на следующем занятии.

	\end{enumerate}
	
	\item Итоги собрания:
	\begin{enumerate}
		
		\item Провод проведен, сервопривод подключен, все работает.
		
		\item Откосы заменены на металлические.
		
        \item Желоб с откидной деталью изготовлены и установлены на робота.
		
	\end{enumerate}
	
	\item Задачи для последующих собраний:
	\begin{enumerate}
		
		\item Реализовать синхронное с опрокидыванием ковша раскрытие направляющей для шаров.
			
	\end{enumerate}
\end{enumerate}
\fillpage
